% ============================================================
% 第2章第7節「特殊関数型」草案
% ============================================================
% このファイルは現行本文には接続していない草案である。
% 2.7 に置くことを想定し、漸化式の形を特殊な恒等式へ翻訳する
% という見方を中心に再構成している。

\section{特殊関数型}

ここまで、漸化式の形を見て、等差型・等比型・階差型・階比型などの
「すでに知っている形」へ変形する方法を学んできた。
しかし、どれだけ式変形をしても、基本パターンに帰着しない漸化式はたくさんある。

例えば、
\[
  a_{n+1}=a_n^2
\]
ならば、$a_n>0$ の範囲では対数を取ることで解ける。
$b_n=\log a_n$ とおけば
\[
  b_{n+1}=2b_n
\]
となり、等比数列に帰着するからである。
一方、
\[
  a_{n+1}=a_n^2+1
\]
では、対数を取っても $+1$ が残る。
これまでの方法が急に働かなくなる。

ところが、係数や定数項が少し変わった
\[
  a_{n+1}=2a_n^2-1
\]
には、別の入口がある。
三角関数の倍角公式
\[
  \cos 2\theta=2\cos^2\theta-1
\]
と右辺が完全に一致している。
そこで $a_n=\cos\theta_n$ と置けば、漸化式は
\[
  \cos\theta_{n+1}=\cos(2\theta_n)
\]
という角度の漸化式へ翻訳できる。

この節で扱うのは、見た目だけでは初等的に解けなさそうでも、
三角関数や多項式、あるいは微分の理論が持つ特別な関係を利用すると、
規則が閉じた形で記述できる漸化式である。
ここではこれを便宜的に「特殊関数型」と呼ぶ。
ただし、特殊関数は問題を解くための魔法の公式ではない。
大切なのは、\emph{どの関数なら、次の項を作る操作を保ったまま表現できるか}を探すことである。

\subsection{特殊関数型を「翻訳」として見る}

漸化式を
\[
  a_{n+1}=F(a_n)
\]
と書く。
これを直接解く代わりに、ある関数 $\varphi$ と別の変化 $G$ を用いて
\[
  a_n=\varphi(t_n),
  \qquad
  t_{n+1}=G(t_n)
\]
と表せないかを考える。

この置き換えが成功するためには、次の関係があればよい。
\[
  F(\varphi(t))=\varphi(G(t)) .
  \tag{2.7.1}
\]
つまり、先に $\varphi$ で数を表してから $F$ を作用させても、
先に $t$ へ $G$ を作用させてから $\varphi$ で数に戻しても、同じ結果になる。

この関係があるとき、$a_1=\varphi(t_1)$ とおけば
\[
  a_n=\varphi\bigl(G^{\,n-1}(t_1)\bigr)
\]
と予想できる。
ここで $G^{\,n-1}$ は $G$ を $n-1$ 回繰り返し作用させるという意味である。
この式は、単なる思いつきではない。
実際、数学的帰納法で確かめられる。

\begin{quote}
  $n=1$ では $a_1=\varphi(t_1)$ である。
  ある $n$ で $a_n=\varphi(G^{\,n-1}(t_1))$ が成り立つとすると、
  \begin{align*}
    a_{n+1}
      &=F(a_n)\\
      &=F\!\left(\varphi\bigl(G^{\,n-1}(t_1)\bigr)\right)\\
      &=\varphi\!\left(G\bigl(G^{\,n-1}(t_1)\bigr)\right)\\
      &=\varphi\bigl(G^{\,n}(t_1)\bigr).
  \end{align*}
  よってすべての $n$ で成立する。
\end{quote}

この節で何度も現れる置き換えは、すべてこの考え方の具体例である。
三角関数を使うときは $\varphi(t)=\cos t$、
指数を使うときは $\varphi(t)=u+u^{-1}$、
ニュートン法を考えるときは「関数の値を $0$ に近づける」という目的に合わせて
新しい数列そのものを設計する。

\subsection{三角関数型：倍角公式を漸化式へ翻訳する}

まず、もっとも分かりやすい例から始めよう。
\[
  a_1=\frac{\sqrt{3}}{2},
  \qquad
  a_{n+1}=2a_n^2-1 .
  \tag{2.7.2}
\]
この漸化式を初項から計算すると、
\[
  \frac{\sqrt{3}}{2},\quad
  \frac12,\quad
  -\frac12,\quad
  -\frac12,\quad\ldots
\]
となる。
しかし、数値を並べるだけでは一般項の仕組みは見えない。

ここで初項に注目する。
\[
  \frac{\sqrt{3}}{2}=\cos\frac{\pi}{6}
\]
である。
また、右辺には
\[
  \cos 2\theta=2\cos^2\theta-1
  \tag{2.7.3}
\]
がそのまま現れている。
したがって、初項の角度を
\[
  \theta_1=\frac{\pi}{6}
\]
とし、角度が毎回2倍になると考える。
すなわち、
\[
  \theta_n=2^{n-1}\theta_1
\]
と置けば、一般項の候補は
\[
  a_n=\cos\left(2^{n-1}\theta_1\right)
  \tag{2.7.4}
\]
である。

ここで、候補を作っただけで終わりにしてはいけない。
漸化式を満たすことを帰納法で確認する。

\begin{quote}
  $n=1$ では、
  \[
    \cos\left(2^{0}\theta_1\right)=\cos\frac{\pi}{6}
    =\frac{\sqrt{3}}{2}=a_1
  \]
  である。
  $a_n=\cos(2^{n-1}\theta_1)$ と仮定すると、倍角公式より
  \begin{align*}
    a_{n+1}
      &=2a_n^2-1\\
      &=2\cos^2\left(2^{n-1}\theta_1\right)-1\\
      &=\cos\left(2^n\theta_1\right).
  \end{align*}
  これは $n+1$ の式である。
\end{quote}

したがって、(2.7.2) の一般項は
\[
  \boxed{
    a_n=\cos\left(\frac{2^{n-1}\pi}{6}\right)
  }
\]
である。

この解法で重要なのは、最初から一般項を当てたことではない。
\begin{enumerate}
  \item 右辺 $2x^2-1$ を見て、倍角公式を思い出す。
  \item 初項を $\cos\theta_1$ の形に直す。
  \item 角度の漸化式 $\theta_{n+1}=2\theta_n$ を考える。
  \item その候補が本当に元の漸化式を満たすことを帰納法で示す。
\end{enumerate}
という流れである。

\subsection{三角関数型の標準形}

初項を一般にして、同じことを行ってみよう。
\[
  a_1=\cos\theta,
  \qquad
  a_{n+1}=2a_n^2-1
  \tag{2.7.5}
\]
とする。
このとき、先ほどの計算から直ちに
\[
  \boxed{a_n=\cos(2^{n-1}\theta)}
  \tag{2.7.6}
\]
が得られる。

初項が $-1\leq a_1\leq1$ なら、少なくとも実数の範囲で
\[
  a_1=\cos\theta
\]
と書ける。
したがって、この範囲の初項を持つ問題では、
\[
  \text{初項を角度に直す}
  \quad\longrightarrow\quad
  \text{倍角公式を使う}
\]
という判断ができる。

ただし、角度 $\theta$ は一意ではない。
$\cos\theta=\cos(-\theta)=\cos(2\pi-\theta)$ だからである。
それでも (2.7.6) はどのような代表角を選んでも同じ $a_n$ を与える。
角度表示は答えそのものというより、漸化式を見通すための座標だと考えるとよい。

\subsection{初項によって姿が変わる：双曲線関数への広がり}

同じ漸化式でも、初項が $[-1,1]$ の外に出ると、
三角関数ではなく双曲線関数が自然になる。
双曲線余弦の倍角公式
\[
  \cosh 2t=2\cosh^2t-1
  \tag{2.7.7}
\]
を見れば、$a_1>1$ のときには
\[
  a_1=\cosh t
\]
と置くのがよいと分かる。
すると全く同じ帰納法により
\[
  a_n=\cosh(2^{n-1}t)
  \tag{2.7.8}
\]
となる。

ここには、初項が解法を選ぶという重要な点がある。
\begin{itemize}
  \item $-1\leq a_1\leq1$ なら $a_1=\cos\theta$ と置く。
  \item $a_1>1$ なら $a_1=\cosh t$ と置く。
  \item $a_1<-1$ なら、まず $a_2=2a_1^2-1>1$ として、$n\geq2$ で双曲線関数を使う。
\end{itemize}

例えば $a_1=3/2$ なら、$t$ を $\cosh t=3/2$ で定めて
\[
  a_n=\cosh(2^{n-1}t)
\]
と書ける。
三角関数の公式だけを暗記していると、初項が $1$ を超えた瞬間に方法が切れてしまう。
しかし本質は、
\[
  F(\varphi(t))=\varphi(2t)
\]
という同じ翻訳関係である。
実数の範囲で使う関数が $\cos$ から $\cosh$ へ変わっただけだ。

\subsection{角度の漸化式として眺める：周期性と不規則さ}

三角関数表示を得ると、数列の振る舞いも角度の運動として見えてくる。
\[
  a_n=\cos\left(2^{n-1}\theta\right)
\]
であるから、角度は毎回2倍される。
ただし、余弦は $2\pi$ 周期なので、角度は
\[
  \theta, 2\theta, 4\theta, 8\theta,ldots
  \pmod{2\pi}
\]
として考えればよい。

例えば $\theta=\pi/6$ のとき、角度を $2\pi$ 周期で見れば
\[
  \frac{\pi}{6},quad
  \frac{\pi}{3},quad
  \frac{2\pi}{3},quad
  \frac{4\pi}{3},quad
  \frac{2\pi}{3},\ldots
\]
となり、数列が途中から $-1/2$ に固定されることが分かる。

一方、$\theta/\pi$ が一般の無理数なら、角度は単純な有限周期を持たない。
同じ漸化式でも、初項の選び方によって、周期的な数列にも、
長く振動し続ける数列にもなる。
「一般項を求める」ことは、数値を一つずつ計算すること以上に、
このような全体の運動を見える形にすることである。

\subsection{対称式型：指数の計算を多項式へ翻訳する}

三角関数を使わず、指数の形から同じ構造を取り出すこともできる。
$u\neq0$ に対して
\[
  x=u+u^{-1}
\]
と置く。
すると
\[
  x^2-2
    =(u+u^{-1})^2-2
    =u^2+u^{-2}
  \tag{2.7.9}
\]
となる。
この一行が、次の漸化式を解く鍵である。
\[
  a_1=u+u^{-1},
  \qquad
  a_{n+1}=a_n^2-2 .
  \tag{2.7.10}
\]

\paragraph{方針}
各項を
\[
  a_n=u^{m_n}+u^{-m_n}
\]
という対称な形に保ちたい。
(2.7.9) に $u^{m_n}$ を代入すれば、指数が2倍されるので、
\[
  m_{n+1}=2m_n,qquad m_1=1
\]
となるはずである。
したがって候補は
\[
  a_n=u^{2^{n-1}}+u^{-2^{n-1}} .
  \tag{2.7.11}
\]

\paragraph{帰納法による確認}
$n=1$ では初期条件そのものである。
また、$a_n=u^{2^{n-1}}+u^{-2^{n-1}}$ と仮定すると、
\begin{align*}
  a_{n+1}
    &=a_n^2-2\\
    &=\left(u^{2^{n-1}}+u^{-2^{n-1}}\right)^2-2\\
    &=u^{2^n}+u^{-2^n}.
\end{align*}
したがって (2.7.11) がすべての $n$ で成立する。

この形が「対称」なのは、$u$ と $u^{-1}$ を入れ替えても値が変わらないからである。
一見すると指数の計算をしているが、$u$ と $u^{-1}$ の対称性があるため、
結果は $x=u+u^{-1}$ の多項式として表せる。

\subsection{三角関数とチェビシェフ多項式の接点}

もし $u=e^{i\theta}$ なら、オイラーの公式から
\[
  u+u^{-1}=2\cos\theta
\]
である。
したがって (2.7.11) は
\[
  a_n=2\cos(2^{n-1}\theta)
\]
と読み替えられる。

さらに、$T_m(x)$ を
\[
  T_m(\cos\theta)=\cos(m\theta)
  \tag{2.7.12}
\]
で定めると、$T_m$ は実際には $x$ の多項式になる。
例えば
\[
  T_0(x)=1,qquad T_1(x)=x,qquad T_2(x)=2x^2-1,qquad
  T_3(x)=4x^3-3x .
\]
倍角公式から
\[
  T_{m+1}(x)=2xT_m(x)-T_{m-1}(x)
  \tag{2.7.13}
\]
という三項漸化式も得られる。

ここで重要なのは、三角関数を別の話題として覚えることではない。
\[
  \text{角度の倍加}
  \quad\longleftrightarrow\quad
  \text{多項式の漸化式}
\]
という翻訳が起きている。
この翻訳をさらに詳しく調べたものがチェビシェフ多項式であり、
付録では漸化式・三角関数・多項式・近似のつながりを改めて扱う。

同じ考え方は3乗にも現れる。
\[
  (u+u^{-1})^3-3(u+u^{-1})=u^3+u^{-3}
\]
だから、
\[
  a_1=u+u^{-1},
  \qquad
  a_{n+1}=a_n^3-3a_n
\]
に対しては
\[
  a_n=u^{3^{n-1}}+u^{-3^{n-1}}
\]
と書ける。
ここでも、漸化式の係数は偶然に決まったのではない。
対称な指数式を次の段階へ送るために、必要な係数として現れている。

\subsection{ニュートン法型：一般項ではなく、近づき方を設計する}

ここまでの例では、特別な関数を使って一般項を明示できた。
しかし、漸化式の価値は一般項が書ける場合だけにあるのではない。
一般項が分からなくても、解へ近づく数列を意図的に設計できる。

$A>0$ の平方根を求めたいとする。
方程式
\[
  f(x)=x^2-A=0
\]
の解を直接求める代わりに、現在の近似値 $a_n$ の近くで $f$ を直線に置き換える。
接線の式は
\[
  y=f(a_n)+f'(a_n)(x-a_n)
\]
である。
この直線が $x$ 軸と交わる点を次の近似値にすれば、
\[
  0=f(a_n)+f'(a_n)(a_{n+1}-a_n)
\]
だから
\[
  a_{n+1}
    =a_n-\frac{f(a_n)}{f'(a_n)}
    =\frac12\left(a_n+\frac{A}{a_n}\right)
  \tag{2.7.14}
\]
となる。
これが平方根に対するニュートン法の漸化式である。

例えば $A=2$、$a_1=2$ とすると、
\[
  a_2=\frac32,qquad
  a_3=\frac{17}{12},qquad
  a_4=\frac{577}{408}
\]
となり、急速に $\sqrt2$ へ近づく。

この収束を、ただの実験で終わらせずに確認しよう。
$a_n>0$ とし、$r=\sqrt A$ と置くと、
\begin{align}
  a_{n+1}-r
    &=\frac{a_n^2+A}{2a_n}-r\\
    &=\frac{a_n^2-2ra_n+r^2}{2a_n}\\
    &=\frac{(a_n-r)^2}{2a_n}.
  \tag{2.7.15}
\end{align}
したがって $a_n\geq r$ なら $a_{n+1}\geq r$ であり、
\[
  a_{n+1}-r
    =\frac{(a_n-r)^2}{2a_n}
    \leq \frac{a_n-r}{2}
\]
となる。
誤差が次の段階でおおよそ二乗されるため、ニュートン法は非常に速い。

ここで得た教訓は、
\[
  \text{「一般項を求める」以外にも、漸化式を設計するという解法がある}
\]
ということである。
特殊な関数が見つからないとき、次の値をどう作ればよいかという視点へ切り替える。

\subsection{似た形を同じ方法で解けるとは限らない}

特殊関数型では、式の見た目が似ていることよりも、
恒等式が正確に一致していることが重要である。
例えば
\[
  2x^2-1
\]
は $\cos 2\theta$ と一致するが、
\[
  2x^2+1
\]
には同じ実三角関数の倍角公式は対応しない。
また、
\[
  x^2-2
\]
には $u+u^{-1}$ の対称式が対応するが、
\[
  x^2+1
\]
はそのままでは同じ対称式を保たない。

もちろん、対応する理論がまだ知られていないという意味で「解けない」と断定してはいけない。
ここで言っているのは、少なくともこの節で導入した置き換えでは閉じない、ということである。
特殊関数型を学ぶ目的は、すべての漸化式に名前を付けることではない。
\begin{quote}
  右辺を見たとき、知っている恒等式のどれと一致するかを調べる。
  一致しなければ、無理に公式を当てはめず、別の見方へ移る。
\end{quote}
という判断力を身につけることである。

\subsection{特殊関数型を解くための手順}

最後に、この節の考え方を実際の解法の手順にまとめる。
\begin{enumerate}
  \item 漸化式の右辺を $F(a_n)$ と見て、どの恒等式の右辺に似ているかを探す。
  \item $a_n=\varphi(t_n)$ と置く候補を決める。
  \item 初項を調べ、$a_1=\varphi(t_1)$ となる $t_1$ を選ぶ。
  \item $F(\varphi(t))=\varphi(G(t))$ が本当に成立するか計算する。
  \item $t_{n+1}=G(t_n)$ を解き、$a_n=\varphi(t_n)$ に戻す。
  \item 最後に数学的帰納法で、初項と漸化式の両方を満たすことを示す。
  \item 三角関数・平方根・分母などの条件を確認する。
\end{enumerate}

特に、最後の帰納法は省略してはいけない。
「倍角公式に似ている」ことは方針を与えるだけであり、
その候補が本当に元の数列であることを保証するのは、初項の確認と漸化式の確認だからである。

漸化式は、項を順番に計算するための規則に見える。
しかし、見方を変えると、ある関数の恒等式を何度も繰り返し適用する装置でもある。
特殊関数型で行っているのは、漸化式を無理に初等的な公式へ押し込むことではない。
その漸化式が本来持っている対称性や倍加構造に合う座標を探し、
そこで規則を読み直すことである。
